% ===========================================================================
%  Programmation Frontend — Lab 4 : Formulaires et React Router
%  ESP/UCAD — 2025–2026
%  Dr. El Hadji Bassirou TOURÉ
%  Compilé avec : xelatex (2–3 passes)
% ===========================================================================
\documentclass[11pt,a4paper]{article}

\usepackage{fontspec}
\setmainfont{Latin Modern Roman}
\setsansfont{Latin Modern Sans}
\setmonofont{Latin Modern Mono}
\usepackage{polyglossia}
\setdefaultlanguage{french}

\usepackage[margin=2.5cm]{geometry}
\usepackage{amsmath,amssymb}
\usepackage{graphicx}
\usepackage{xcolor}
\usepackage{booktabs}
\usepackage{hyperref}
\usepackage{tcolorbox}
\tcbuselibrary{theorems,breakable,skins}
\usepackage{enumitem}
\usepackage{fancyhdr}
\usepackage{array}
\usepackage{pifont}

\newcommand{\cmark}{\ding{51}}
\newcommand{\xmark}{\ding{55}}

\hypersetup{colorlinks=true,linkcolor=blue!60!black,urlcolor=blue!60!black,citecolor=blue!60!black}

\newtcolorbox{objectifbox}[1][]{colback=blue!5,colframe=blue!50,fonttitle=\bfseries,title=#1,breakable}
\newtcolorbox{infobox}[1][]{colback=green!5,colframe=green!50!black,fonttitle=\bfseries,title=#1,breakable}
\newtcolorbox{attentionbox}[1][]{colback=orange!5,colframe=orange!70,fonttitle=\bfseries,title=#1,breakable}
\newtcolorbox{retenirbox}[1][]{colback=yellow!10,colframe=yellow!50!black,fonttitle=\bfseries,title=#1,breakable}
\newtcolorbox{prereqbox}[1][]{colback=cyan!5,colframe=cyan!50!black,fonttitle=\bfseries,title=#1,breakable}
\newtcolorbox{projetbox}[1][]{colback=teal!5,colframe=teal!50!black,fonttitle=\bfseries,title=#1,breakable}

\usepackage{listings}
\lstdefinelanguage{JavaScript}{keywords={break,case,catch,continue,debugger,default,delete,do,else,finally,for,function,if,in,instanceof,new,return,switch,this,throw,try,typeof,var,void,while,with,const,let,class,export,extends,import,super,yield,async,await,of,from,true,false,null,undefined},morecomment=[l]{//},morecomment=[s]{/*}{*/},morestring=[b]',morestring=[b]",morestring=[b]`,sensitive=true}
\lstdefinelanguage{CSS}{keywords={color,background,border,padding,margin,font,display,flex,align,justify,width,height,gap,cursor,none,solid,center,column,row,wrap,pointer,radius,size,weight,family,items,content,direction,transform,transition,hover,box,shadow,max,min,text,decoration,opacity,position,sticky,top,letter,spacing,uppercase},morecomment=[s]{/*}{*/},morestring=[b]',morestring=[b]",sensitive=true}
\lstset{basicstyle=\ttfamily\small,backgroundcolor=\color{gray!8},frame=single,framerule=0.5pt,rulecolor=\color{gray!40},breaklines=true,breakatwhitespace=true,showstringspaces=false,tabsize=2,xleftmargin=4pt,xrightmargin=4pt,aboveskip=8pt,belowskip=8pt,language=JavaScript,keywordstyle=\color{blue!70!black}\bfseries,commentstyle=\color{green!50!black}\itshape,stringstyle=\color{red!60!black},literate={é}{{\'e}}1 {è}{{\`e}}1 {ê}{{\^e}}1 {à}{{\`a}}1 {ù}{{\`u}}1 {ô}{{\^o}}1 {î}{{\^i}}1 {ç}{{\c{c}}}1 {É}{{\'E}}1 {→}{{$\rightarrow$}}1}
\lstdefinestyle{terminal}{language={},keywordstyle={},commentstyle={},stringstyle={},basicstyle=\ttfamily\small,backgroundcolor=\color{black!5}}
\lstdefinestyle{css}{language=CSS,keywordstyle=\color{blue!70!black}\bfseries,commentstyle=\color{green!50!black}\itshape,stringstyle=\color{red!60!black}}

\pagestyle{fancy}\fancyhf{}
\fancyhead[L]{\small Prog. Frontend --- Lab 4 : Formulaires et React Router}
\fancyhead[R]{\small ESP/UCAD}
\fancyfoot[C]{\thepage}
\fancyfoot[L]{\small Dr. El Hadji Bassirou TOURÉ}
\fancyfoot[R]{\small 2025--2026}
\renewcommand{\headrulewidth}{0.4pt}
\renewcommand{\footrulewidth}{0.4pt}

\title{\vspace{-1cm}\textbf{\LARGE Programmation Frontend}\\[0.8cm]\textbf{\huge Lab 4 --- Formulaires et React Router}\\[0.3cm]\textbf{\Large Multi-pages, navigation et saisie de données}\\[0.5cm]\large ESP/UCAD --- 2025--2026}
\author{Dr. El Hadji Bassirou TOURÉ\\Département Génie Informatique\\Université Cheikh Anta Diop de Dakar}
\date{}

\begin{document}
\maketitle
\thispagestyle{empty}
\vspace{0.5cm}

\begin{prereqbox}[Informations pratiques]
\begin{center}
\begin{tabular}{lp{10cm}}
\toprule
\textbf{Durée estimée} & 2h (une séance) \\
\textbf{Prérequis} & Avoir terminé le Lab 3 (\texttt{useEffect}, loading/erreur, vue détail). \\
\textbf{Fichiers de départ} & Le projet \texttt{annuaire-react/} du Lab 3 \\
\textbf{Livrables} & SPA avec 4 pages (accueil, liste, détail, ajout), navigation et formulaire \\
\textbf{Validation} & Chaque vue a sa propre URL. Le formulaire valide et envoie les données. \\
\bottomrule
\end{tabular}
\end{center}
\end{prereqbox}

\begin{objectifbox}[Objectifs du lab]
À l'issue de ce lab, chaque étudiant sera capable de :
\begin{enumerate}[nosep]
\item Installer et configurer \textbf{React Router} pour créer une application multi-pages
\item Créer des \textbf{routes}, des \textbf{liens} (\texttt{Link}) et une \textbf{barre de navigation}
\item Lire les \textbf{paramètres d'URL} avec \texttt{useParams()}
\item Créer un \textbf{formulaire contrôlé} avec \texttt{useState}
\item \textbf{Valider} les champs et afficher les erreurs
\item Envoyer des données avec \texttt{fetch} en méthode \texttt{POST}
\end{enumerate}
\end{objectifbox}

\tableofcontents
\newpage

% ===========================================================================
\section{Point de départ --- Le problème du Lab 3}
% ===========================================================================

Au Lab 3, la navigation entre la liste et le détail est gérée par un état \texttt{selectedUserId}. Ça fonctionne, mais :

\begin{itemize}[nosep]
\item L'\textbf{URL ne change pas}. Liste ou détail, l'adresse reste \texttt{localhost:5173}.
\item Le \textbf{bouton retour} du navigateur ne fonctionne pas.
\item On ne peut pas \textbf{partager un lien} vers un utilisateur précis.
\item On ne peut pas \textbf{ouvrir un profil dans un nouvel onglet}.
\end{itemize}

React Router résout tout cela : chaque vue obtient sa propre URL.

\begin{objectifbox}[Étape 0 --- Installation de React Router]
Dans le terminal, à la racine du projet :

\begin{lstlisting}[style=terminal]
cd annuaire-react
npm install react-router-dom
\end{lstlisting}

Vérifiez que la ligne \texttt{"react-router-dom"} apparaît dans \texttt{package.json} sous \texttt{dependencies}.
\end{objectifbox}

% ===========================================================================
\section{Partie 1 --- Mettre en place le routeur}
\label{sec:routeur}
% ===========================================================================

\begin{objectifbox}[Objectif --- 10 minutes]
Configurer React Router et définir les routes de l'application.
\end{objectifbox}

\subsection{Envelopper l'application dans \texttt{BrowserRouter}}

Modifiez \texttt{src/main.jsx} :

\begin{lstlisting}
import { StrictMode } from "react";
import { createRoot } from "react-dom/client";
import { BrowserRouter } from "react-router-dom";
import App from "./App";

createRoot(document.getElementById("root")).render(
  <StrictMode>
    <BrowserRouter>
      <App />
    </BrowserRouter>
  </StrictMode>
);
\end{lstlisting}

\texttt{BrowserRouter} active le système de routing dans toute l'application. Tous les composants enfants pourront utiliser \texttt{Link}, \texttt{useParams}, etc.

\subsection{Créer les dossiers}

\begin{lstlisting}[style=terminal]
mkdir -p src/pages
\end{lstlisting}

On sépare les \textbf{pages} (composants liés à une route) des \textbf{composants} réutilisables.

\subsection{Définir les routes dans \texttt{App.jsx}}

Remplacez \textbf{tout le contenu} de \texttt{src/App.jsx} :

\begin{lstlisting}
import { Routes, Route } from "react-router-dom";
import Navbar from "./components/Navbar";
import Home from "./pages/Home";
import UserList from "./pages/UserList";
import UserDetail from "./pages/UserDetail";
import AddUser from "./pages/AddUser";
import styles from "./App.module.css";

const App = () => {
  return (
    <div className={styles.app}>
      <Navbar />
      <main className={styles.main}>
        <Routes>
          <Route path="/" element={<Home />} />
          <Route path="/users"
                 element={<UserList />} />
          <Route path="/users/:id"
                 element={<UserDetail />} />
          <Route path="/users/add"
                 element={<AddUser />} />
        </Routes>
      </main>
    </div>
  );
};

export default App;
\end{lstlisting}

\begin{retenirbox}[Les routes de l'annuaire]
\begin{center}
\begin{tabular}{llp{5.5cm}}
\toprule
\textbf{URL} & \textbf{Page} & \textbf{Description} \\
\midrule
\texttt{/} & \texttt{Home} & Page d'accueil \\
\texttt{/users} & \texttt{UserList} & Liste avec recherche, tri et favoris \\
\texttt{/users/:id} & \texttt{UserDetail} & Détail d'un utilisateur (\texttt{:id} est un paramètre) \\
\texttt{/users/add} & \texttt{AddUser} & Formulaire d'ajout \\
\bottomrule
\end{tabular}
\end{center}
\end{retenirbox}

Remplacez \texttt{src/App.module.css} :

\begin{lstlisting}[style=css]
.app {
  min-height: 100vh;
  background: #f5f6fa;
}

.main {
  max-width: 700px;
  margin: 0 auto;
  padding: 2rem 1rem;
}
\end{lstlisting}

Ajoutez aussi un reset CSS global dans \texttt{index.html}, dans la balise \texttt{<head>} :

\begin{lstlisting}[language=HTML]
<style>
  *, *::before, *::after {
    box-sizing: border-box; margin: 0; padding: 0;
  }
  body {
    font-family: -apple-system, BlinkMacSystemFont,
      "Segoe UI", Roboto, sans-serif;
    color: #2d3436;
    background: #f5f6fa;
    line-height: 1.6;
  }
  a { text-decoration: none; color: inherit; }
</style>
\end{lstlisting}

\begin{infobox}[App.jsx a changé de rôle]
Au Lab 3, \texttt{App.jsx} contenait toute la logique (états, fetch, filtrage). Maintenant, il ne fait que \textbf{deux choses} : afficher la \texttt{Navbar} et définir les \texttt{Routes}. Chaque page gère sa propre logique. Le code est mieux organisé.
\end{infobox}

% ===========================================================================
\section{Partie 2 --- Navigation avec Navbar et Link}
\label{sec:navbar}
% ===========================================================================

\begin{objectifbox}[Objectif --- 10 minutes]
Créer une barre de navigation qui fonctionne sans rechargement de page.
\end{objectifbox}

Créez \texttt{src/components/Navbar.jsx} :

\begin{lstlisting}
import { NavLink } from "react-router-dom";
import styles from "./Navbar.module.css";

const Navbar = () => {
  return (
    <nav className={styles.nav}>
      <div className={styles.container}>
        <NavLink to="/" className={styles.logo}>
          Annuaire
        </NavLink>
        <div className={styles.links}>
          <NavLink to="/users"
            className={({ isActive }) =>
              isActive
                ? `${styles.link} ${styles.active}`
                : styles.link
            }>
            Liste
          </NavLink>
          <NavLink to="/users/add"
            className={({ isActive }) =>
              isActive
                ? `${styles.link} ${styles.active}`
                : styles.link
            }>
            + Ajouter
          </NavLink>
        </div>
      </div>
    </nav>
  );
};

export default Navbar;
\end{lstlisting}

\begin{attentionbox}[\texttt{Link} vs \texttt{NavLink} vs \texttt{<a>}]
\begin{center}
\begin{tabular}{lp{8cm}}
\toprule
\textbf{Élément} & \textbf{Usage} \\
\midrule
\texttt{<a href="...">} & Lien HTML classique. \textbf{Recharge} toute la page. Ne jamais utiliser pour la navigation interne. \\
\texttt{<Link to="...">} & Lien React Router. Change l'URL \textbf{sans} recharger la page. \\
\texttt{<NavLink to="...">} & Comme \texttt{Link}, mais fournit en plus \texttt{isActive} pour styliser le lien de la page courante. \\
\bottomrule
\end{tabular}
\end{center}
\end{attentionbox}

Créez \texttt{src/components/Navbar.module.css} :

\begin{lstlisting}[style=css]
.nav {
  background: #2d3436;
  padding: 0 1rem;
  position: sticky;
  top: 0;
  z-index: 10;
}

.container {
  max-width: 700px;
  margin: 0 auto;
  display: flex;
  justify-content: space-between;
  align-items: center;
  height: 56px;
}

.logo {
  font-size: 1.25rem;
  font-weight: 700;
  color: #fff;
}

.links { display: flex; gap: 0.25rem; }

.link {
  color: #b2bec3;
  padding: 0.4rem 0.9rem;
  border-radius: 6px;
  font-size: 0.9rem;
  transition: background 0.15s, color 0.15s;
}

.link:hover {
  color: #fff;
  background: rgba(255, 255, 255, 0.1);
}

.active {
  color: #fff;
  background: rgba(255, 255, 255, 0.15);
}
\end{lstlisting}

% ===========================================================================
\section{Partie 3 --- Les pages}
\label{sec:pages}
% ===========================================================================

\begin{objectifbox}[Objectif --- 15 minutes]
Créer les pages Home, UserList et déplacer UserDetail.
\end{objectifbox}

\subsection{Page d'accueil}

Créez \texttt{src/pages/Home.jsx} :

\begin{lstlisting}
import { Link } from "react-router-dom";
import styles from "./Home.module.css";

const Home = () => {
  return (
    <div className={styles.home}>
      <h1 className={styles.titre}>
        Annuaire des utilisateurs
      </h1>
      <p className={styles.description}>
        Consultez, recherchez et ajoutez des utilisateurs.
      </p>
      <div className={styles.actions}>
        <Link to="/users" className={styles.bouton}>
          Voir la liste
        </Link>
        <Link to="/users/add"
              className={styles.boutonSecondaire}>
          Ajouter un utilisateur
        </Link>
      </div>
    </div>
  );
};

export default Home;
\end{lstlisting}

Créez \texttt{src/pages/Home.module.css} :

\begin{lstlisting}[style=css]
.home { text-align: center; padding: 3rem 1rem; }

.titre {
  font-size: 2rem; color: #2d3436;
  margin-bottom: 1rem;
}

.description {
  color: #636e72; font-size: 1.05rem;
  max-width: 450px; margin: 0 auto 2rem auto;
}

.actions {
  display: flex; gap: 1rem;
  justify-content: center; flex-wrap: wrap;
}

.bouton {
  padding: 0.7rem 1.8rem; background: #0984e3;
  color: #fff; border-radius: 8px;
  font-weight: 500; transition: background 0.15s;
}
.bouton:hover { background: #0770c2; }

.boutonSecondaire {
  padding: 0.7rem 1.8rem; background: #fff;
  color: #0984e3; border: 1px solid #0984e3;
  border-radius: 8px; font-weight: 500;
}
.boutonSecondaire:hover { background: #ebf5fb; }
\end{lstlisting}

\subsection{Page de liste}

Créez \texttt{src/pages/UserList.jsx}. Cette page reprend la logique qui était dans \texttt{App.jsx} au Lab 3 :

\begin{lstlisting}
import { useState, useEffect } from "react";
import UserCard from "../components/UserCard";
import SearchBar from "../components/SearchBar";
import styles from "./UserList.module.css";

const UserList = () => {
  const [users, setUsers] = useState([]);
  const [chargement, setChargement] = useState(true);
  const [erreur, setErreur] = useState(null);
  const [recherche, setRecherche] = useState("");
  const [triAscendant, setTriAscendant] = useState(true);
  const [favoris, setFavoris] = useState([]);

  useEffect(() => {
    const charger = async () => {
      try {
        const response = await fetch(
          "https://jsonplaceholder.typicode.com/users"
        );
        if (!response.ok)
          throw new Error("Erreur serveur");
        const data = await response.json();
        setUsers(data);
      } catch (error) {
        setErreur(error.message);
      }
      setChargement(false);
    };
    charger();
  }, []);

  const toggleFavori = (userId) => {
    if (favoris.includes(userId)) {
      setFavoris(
        favoris.filter((id) => id !== userId)
      );
    } else {
      setFavoris([...favoris, userId]);
    }
  };

  if (chargement)
    return <p>Chargement en cours...</p>;
  if (erreur)
    return <p className={styles.erreur}>
      Erreur : {erreur}</p>;

  const usersFiltres = users
    .filter((u) =>
      u.name.toLowerCase()
        .includes(recherche.toLowerCase())
    )
    .sort((a, b) =>
      triAscendant
        ? a.name.localeCompare(b.name)
        : b.name.localeCompare(a.name)
    );

  return (
    <div>
      <h1 className={styles.titre}>
        Liste des utilisateurs
      </h1>
      <SearchBar recherche={recherche}
                 onRecherche={setRecherche} />
      <div className={styles.barre}>
        <span className={styles.compteur}>
          {usersFiltres.length} utilisateur(s)
        </span>
        {favoris.length > 0 && (
          <span className={styles.favoris}>
            {favoris.length} favori(s)
          </span>
        )}
        <button className={styles.boutonTri}
                onClick={() =>
                  setTriAscendant(!triAscendant)
                }>
          {triAscendant ? "A → Z" : "Z → A"}
        </button>
      </div>
      {usersFiltres.map((user) => (
        <UserCard key={user.id} user={user}
          estFavori={favoris.includes(user.id)}
          onToggleFavori={toggleFavori} />
      ))}
    </div>
  );
};

export default UserList;
\end{lstlisting}

\subsection{Modifier \texttt{UserCard} pour utiliser \texttt{Link}}

Au Lab 3, les cartes utilisaient \texttt{onClick} pour naviguer par état. Maintenant, on utilise \texttt{Link} pour naviguer par URL. Modifiez \texttt{src/components/UserCard.jsx} :

\begin{lstlisting}
import { Link } from "react-router-dom";
import styles from "./UserCard.module.css";

const UserCard = ({
  user, estFavori, onToggleFavori
}) => {
  return (
    <div className={styles.card}>
      <div className={styles.header}>
        <Link to={`/users/${user.id}`}
              className={styles.lien}>
          <h3 className={styles.nom}>{user.name}</h3>
        </Link>
        <button className={styles.etoile}
                onClick={() =>
                  onToggleFavori(user.id)
                }>
          {estFavori ? "\u2605" : "\u2606"}
        </button>
      </div>
      <p className={styles.info}>
        Email : {user.email}
      </p>
      <p className={styles.info}>
        Entreprise : {user.company.name}
      </p>
      <Link to={`/users/${user.id}`}
            className={styles.voir}>
        Voir le profil →
      </Link>
    </div>
  );
};

export default UserCard;
\end{lstlisting}

\begin{infobox}[Plus besoin de \texttt{e.stopPropagation()}]
Au Lab 3, le clic sur l'étoile déclenchait aussi le \texttt{onClick} de la carte. On avait besoin de \texttt{e.stopPropagation()}. Maintenant, la navigation se fait par \texttt{Link} (un lien, pas un \texttt{onClick} sur la carte entière). Le bouton étoile n'entre plus en conflit.
\end{infobox}

\subsection{Page de détail avec \texttt{useParams}}

Créez \texttt{src/pages/UserDetail.jsx}. Le composant lit l'identifiant \textbf{depuis l'URL} au lieu de le recevoir en prop :

\begin{lstlisting}
import { useState, useEffect } from "react";
import { useParams, Link } from "react-router-dom";
import styles from "./UserDetail.module.css";

const UserDetail = () => {
  const { id } = useParams();
  const [user, setUser] = useState(null);
  const [chargement, setChargement] = useState(true);
  const [erreur, setErreur] = useState(null);

  useEffect(() => {
    const charger = async () => {
      setChargement(true);
      try {
        const response = await fetch(
          `https://jsonplaceholder.typicode.com/users/${id}`
        );
        if (!response.ok)
          throw new Error("Utilisateur introuvable");
        const data = await response.json();
        setUser(data);
      } catch (error) {
        setErreur(error.message);
      }
      setChargement(false);
    };
    charger();
  }, [id]);

  if (chargement) return <p>Chargement...</p>;

  if (erreur) {
    return (
      <div>
        <p>{erreur}</p>
        <Link to="/users">Retour à la liste</Link>
      </div>
    );
  }

  return (
    <div>
      <Link to="/users" className={styles.retour}>
        ← Retour à la liste
      </Link>
      <div className={styles.profil}>
        <div className={styles.avatar}>
          {user.name.charAt(0)}
        </div>
        <h1 className={styles.nom}>{user.name}</h1>
        <p className={styles.username}>
          @{user.username}
        </p>
      </div>
      <div className={styles.carte}>
        <h2 className={styles.section}>Coordonnées</h2>
        <div className={styles.champ}>
          <span className={styles.label}>Email</span>
          <span>{user.email}</span>
        </div>
        <div className={styles.champ}>
          <span className={styles.label}>Tél.</span>
          <span>{user.phone}</span>
        </div>
        <div className={styles.champ}>
          <span className={styles.label}>Site</span>
          <span>{user.website}</span>
        </div>
      </div>
      <div className={styles.carte}>
        <h2 className={styles.section}>Entreprise</h2>
        <div className={styles.champ}>
          <span className={styles.label}>Nom</span>
          <span>{user.company.name}</span>
        </div>
        <div className={styles.champ}>
          <span className={styles.label}>Ville</span>
          <span>{user.address.city}</span>
        </div>
      </div>
    </div>
  );
};

export default UserDetail;
\end{lstlisting}

\begin{retenirbox}[Lab 3 vs Lab 4 : la page de détail]
\begin{center}
\begin{tabular}{lp{4.5cm}p{4.5cm}}
\toprule
& \textbf{Lab 3} & \textbf{Lab 4} \\
\midrule
\textbf{L'identifiant vient de} & Une prop \texttt{userId} (passée par App) & L'URL via \texttt{useParams()} \\
\textbf{L'URL} & Ne change pas & \texttt{/users/3} \\
\textbf{Bouton retour navigateur} & Ne fonctionne pas & Fonctionne \\
\textbf{Partage de lien} & Impossible & Possible \\
\bottomrule
\end{tabular}
\end{center}
\end{retenirbox}

Supprimez l'ancien \texttt{src/components/UserDetail.jsx} et son CSS Module --- il est remplacé par la page. Les fichiers CSS des pages (\texttt{UserDetail.module.css}, \texttt{UserList.module.css}) sont fournis dans le script de setup.

% ===========================================================================
\section{Partie 4 --- Formulaire contrôlé}
\label{sec:formulaire}
% ===========================================================================

\begin{objectifbox}[Objectif --- 20 minutes]
Créer un formulaire pour ajouter un utilisateur. Comprendre les inputs contrôlés.
\end{objectifbox}

\subsection{Input contrôlé vs non contrôlé}

Un \textbf{input contrôlé} est un champ dont la valeur est stockée dans l'état React. À chaque frappe, l'état se met à jour, et React re-rend le composant avec la nouvelle valeur.

\begin{center}
\begin{tabular}{lp{5cm}p{5cm}}
\toprule
& \textbf{Non contrôlé} & \textbf{Contrôlé} \\
\midrule
\textbf{Valeur stockée dans} & Le DOM (le navigateur) & L'état React \\
\textbf{Lecture} & \texttt{inputRef.current.value} & \texttt{nom} (variable d'état) \\
\textbf{Avantage} & Simple pour les cas basiques & Total contrôle (validation, formatage) \\
\bottomrule
\end{tabular}
\end{center}

En React, on utilise \textbf{toujours} des inputs contrôlés. Le principe : \texttt{value=\{état\}} + \texttt{onChange=\{fn\}}.

\subsection{Créer \texttt{AddUser.jsx}}

Créez \texttt{src/pages/AddUser.jsx} :

\begin{lstlisting}
import { useState } from "react";
import { useNavigate } from "react-router-dom";
import styles from "./AddUser.module.css";

const AddUser = () => {
  const navigate = useNavigate();

  const [form, setForm] = useState({
    name: "",
    email: "",
    phone: "",
    company: "",
  });
  const [erreurs, setErreurs] = useState({});
  const [enCours, setEnCours] = useState(false);
  const [soumis, setSoumis] = useState(false);

  const handleChange = (e) => {
    const { name, value } = e.target;
    setForm({ ...form, [name]: value });
    if (erreurs[name]) {
      setErreurs({ ...erreurs, [name]: null });
    }
  };

  const valider = () => {
    const err = {};
    if (!form.name.trim())
      err.name = "Le nom est obligatoire";
    if (!form.email.trim())
      err.email = "L'email est obligatoire";
    else if (!form.email.includes("@"))
      err.email = "L'email doit contenir un @";
    if (!form.phone.trim())
      err.phone = "Le téléphone est obligatoire";
    return err;
  };

  const handleSubmit = async (e) => {
    e.preventDefault();

    const err = valider();
    if (Object.keys(err).length > 0) {
      setErreurs(err);
      return;
    }

    setEnCours(true);
    try {
      const response = await fetch(
        "https://jsonplaceholder.typicode.com/users",
        {
          method: "POST",
          headers: {
            "Content-Type": "application/json",
          },
          body: JSON.stringify({
            name: form.name,
            email: form.email,
            phone: form.phone,
            company: {
              name: form.company || "Non renseigné",
            },
          }),
        }
      );
      if (!response.ok)
        throw new Error("Erreur serveur");

      const data = await response.json();
      console.log("Utilisateur créé :", data);
      setSoumis(true);
      setTimeout(() => navigate("/users"), 2000);
    } catch (error) {
      setErreurs({ submit: error.message });
    }
    setEnCours(false);
  };

  if (soumis) {
    return (
      <div className={styles.succes}>
        <h2>Utilisateur ajouté</h2>
        <p>Redirection vers la liste...</p>
      </div>
    );
  }

  return (
    <div>
      <h1 className={styles.titre}>
        Ajouter un utilisateur
      </h1>
      <div className={styles.carte}>
        <form onSubmit={handleSubmit}>
          <div className={styles.groupe}>
            <label className={styles.label}
                   htmlFor="name">
              Nom complet *
            </label>
            <input className={styles.input}
              type="text" id="name" name="name"
              placeholder="Fatou Ndiaye"
              value={form.name}
              onChange={handleChange} />
            {erreurs.name && (
              <p className={styles.erreur}>
                {erreurs.name}
              </p>
            )}
          </div>

          <div className={styles.groupe}>
            <label className={styles.label}
                   htmlFor="email">
              Email *
            </label>
            <input className={styles.input}
              type="text" id="email" name="email"
              placeholder="fatou.ndiaye@ucad.edu.sn"
              value={form.email}
              onChange={handleChange} />
            {erreurs.email && (
              <p className={styles.erreur}>
                {erreurs.email}
              </p>
            )}
          </div>

          <div className={styles.groupe}>
            <label className={styles.label}
                   htmlFor="phone">
              Téléphone *
            </label>
            <input className={styles.input}
              type="text" id="phone" name="phone"
              placeholder="77 123 45 67"
              value={form.phone}
              onChange={handleChange} />
            {erreurs.phone && (
              <p className={styles.erreur}>
                {erreurs.phone}
              </p>
            )}
          </div>

          <div className={styles.groupe}>
            <label className={styles.label}
                   htmlFor="company">
              Entreprise
            </label>
            <input className={styles.input}
              type="text" id="company" name="company"
              placeholder="ESP/UCAD"
              value={form.company}
              onChange={handleChange} />
          </div>

          {erreurs.submit && (
            <p className={styles.erreurGlobale}>
              {erreurs.submit}
            </p>
          )}

          <button type="submit"
                  className={styles.bouton}
                  disabled={enCours}>
            {enCours ? "Envoi..." : "Ajouter"}
          </button>
        </form>
      </div>
    </div>
  );
};

export default AddUser;
\end{lstlisting}

\begin{retenirbox}[Décoder le formulaire]
\begin{center}
\begin{tabular}{lp{8cm}}
\toprule
\textbf{Élément} & \textbf{Rôle} \\
\midrule
\texttt{value=\{form.name\}} & L'input affiche la valeur de l'état (input contrôlé) \\
\texttt{onChange=\{handleChange\}} & À chaque frappe, met à jour l'état \\
\texttt{e.preventDefault()} & Empêche le rechargement de la page à la soumission \\
\texttt{name="name"} & L'attribut \texttt{name} permet à \texttt{handleChange} de savoir quel champ mettre à jour \\
\texttt{[name]: value} & Syntaxe ES6 : clé dynamique dans un objet \\
\texttt{useNavigate()} & Permet de naviguer par code (après soumission) \\
\texttt{method: "POST"} & Envoie les données au serveur (au lieu de les lire) \\
\bottomrule
\end{tabular}
\end{center}
\end{retenirbox}

\begin{attentionbox}[JSONPlaceholder simule le POST]
L'API JSONPlaceholder accepte les requêtes POST et retourne un objet avec un \texttt{id} généré. Mais les données ne sont \textbf{pas réellement enregistrées}. Si vous rechargez la liste, le nouvel utilisateur n'apparaît pas. C'est normal --- c'est une API de test. Au Lab 6, on se connectera à un vrai backend Django.
\end{attentionbox}

Les fichiers CSS pour les pages (\texttt{AddUser.module.css}, \texttt{UserDetail.module.css}, \texttt{UserList.module.css}) sont fournis dans le script de setup. Ils suivent la même logique que les labs précédents : un fichier \texttt{.module.css} par composant.

% ===========================================================================
\section{Bilan --- Structure du projet}
% ===========================================================================

\begin{lstlisting}[style=terminal]
annuaire-react/
  src/
    App.jsx               <-- Routes seulement
    App.module.css
    components/
      Navbar.jsx          <-- nouveau
      Navbar.module.css   <-- nouveau
      UserCard.jsx        <-- modifié (Link)
      UserCard.module.css
      SearchBar.jsx
      SearchBar.module.css
    pages/                <-- nouveau dossier
      Home.jsx
      Home.module.css
      UserList.jsx        <-- logique déplacée ici
      UserList.module.css
      UserDetail.jsx      <-- useParams
      UserDetail.module.css
      AddUser.jsx         <-- nouveau
      AddUser.module.css
  index.html
  package.json
\end{lstlisting}

\begin{retenirbox}[Avant / après : Lab 3 vs Lab 4]
\begin{center}
\begin{tabular}{lp{4.5cm}p{4.5cm}}
\toprule
& \textbf{Lab 3} & \textbf{Lab 4} \\
\midrule
\textbf{Navigation} & État \texttt{selectedUserId} & URL (\texttt{/users}, \texttt{/users/3}) \\
\textbf{Bouton retour} & Ne fonctionne pas & Fonctionne \\
\textbf{Partage de lien} & Impossible & \texttt{/users/3} \\
\textbf{Ajout} & Aucun & Formulaire contrôlé + POST \\
\textbf{Lien actif} & Aucun & NavLink avec style actif \\
\textbf{App.jsx} & Toute la logique & Routes seulement \\
\bottomrule
\end{tabular}
\end{center}
\end{retenirbox}

% ===========================================================================
\section{Aide-mémoire}
% ===========================================================================

\begin{center}
\begin{tabular}{lp{7.5cm}}
\toprule
\textbf{Concept} & \textbf{Syntaxe} \\
\midrule
Installer & \texttt{npm install react-router-dom} \\
Envelopper & \texttt{<BrowserRouter><App /></BrowserRouter>} \\
Définir une route & \texttt{<Route path="/x" element=\{<Page />\} />} \\
Paramètre d'URL & \texttt{path="/users/:id"} \\
Lire le paramètre & \texttt{const \{ id \} = useParams();} \\
Lien sans rechargement & \texttt{<Link to="/users">...</Link>} \\
Lien avec style actif & \texttt{<NavLink to="/users" className=\{...\}>} \\
Navigation par code & \texttt{const navigate = useNavigate(); navigate("/users");} \\
\midrule
Input contrôlé & \texttt{<input value=\{état\} onChange=\{fn\} />} \\
Empêcher le rechargement & \texttt{e.preventDefault()} \\
Clé dynamique & \texttt{setForm(\{ ...form, [name]: value \})} \\
POST avec fetch & \texttt{fetch(url, \{ method: "POST", body: JSON.stringify(data) \})} \\
\bottomrule
\end{tabular}
\end{center}

\newpage

% ===========================================================================
\section{Résumé et conclusion}
% ===========================================================================

Dans ce lab, vous avez :
\begin{enumerate}[nosep]
\item Installé \textbf{React Router} et défini 4 routes avec leurs URL.
\item Créé une \textbf{Navbar} avec des \texttt{NavLink} qui indiquent la page active.
\item Utilisé \texttt{useParams()} pour lire l'identifiant de l'utilisateur depuis l'URL.
\item Créé un \textbf{formulaire contrôlé} avec validation et envoi en POST.
\item Utilisé \texttt{useNavigate()} pour rediriger après la soumission.
\end{enumerate}

\begin{center}
\begin{tabular}{lcp{7cm}}
\toprule
\textbf{Étape} & \textbf{Statut} & \textbf{Ce que vous savez faire} \\
\midrule
\textbf{DOM et événements} & \cmark{} & Sélecteurs, événements \\
\textbf{Fonctions JS} & \cmark{} & Déclaration, expression, fléchée \\
\textbf{Asynchrone} & \cmark{} & \texttt{fetch}, \texttt{async/await} \\
\textbf{React --- les bases} & \cmark{} & Composant, JSX, Vite \\
\textbf{Composants et props} & \cmark{} & Découpage, props, CSS Modules \\
\textbf{État React} & \cmark{} & \texttt{useState}, re-rendu, immutabilité \\
\textbf{useEffect} & \cmark{} & Montage, dépendances, loading/erreur \\
\textbf{Router et formulaires} & \cmark{} & Routes, Link, useParams, inputs contrôlés, POST \\
\textbf{Next.js} & $\to$ Lab 5 & App Router, server vs client \\
\textbf{API Django} & $\to$ Lab 6 & CORS, auth, CRUD \\
\bottomrule
\end{tabular}
\end{center}

\begin{projetbox}[Et ensuite ?]
L'annuaire est une SPA complète avec 4 pages, de la navigation, un formulaire et des appels API. Mais React seul a des limites : pas de SEO (le contenu est généré côté client), un premier chargement lent (tout le JS est téléchargé d'un coup), et un routing configuré manuellement.

Dans le \textbf{Lab 5}, vous migrerez l'annuaire vers \textbf{Next.js} : routing par fichiers, layouts, composants serveur, et \texttt{loading.js} automatique.
\end{projetbox}

\section*{Exercice bonus}

\begin{infobox}[Bonus : page 404]
Ajoutez une route «~attrape-tout~» pour les URL invalides :
\begin{lstlisting}
<Route path="*" element={<NotFound />} />
\end{lstlisting}
Créez \texttt{src/pages/NotFound.jsx} avec un message «~Page introuvable~» et un lien vers l'accueil. Testez en tapant \texttt{/xyz} dans l'URL.
\end{infobox}

\vspace{1cm}

\begin{center}
\fbox{
\begin{minipage}{0.85\textwidth}
\centering
\vspace{0.5cm}
\textbf{\Large Lab 4 --- Formulaires et React Router --- Terminé \cmark}\\[0.3cm]
\normalsize
Vous savez créer une application multi-pages avec\\
des routes, de la navigation, et des formulaires contrôlés.\\[0.2cm]
Prochaine étape : Next.js et le rendu côté serveur.\\[0.3cm]
\textit{Chaque page a son URL. Chaque URL a sa page.}\\[0.5cm]
\end{minipage}
}
\end{center}

\end{document}
